\chapter{Γ 函数.}

对序列 \((u_{n})\) 用一个依赖于实参数 \(\lambda\) 且等于 \(u_{n}\)（当 \(\lambda = n\) 时）的积分的值来进行“插值”的想法，可追溯到沃利斯（参见第 187 页以下）。正是这个想法在 1730 年（{[}108 a{]}，(1)，第 XIV 卷，第 1–24 页）成为引导欧拉对阶乘序列进行插值的主要线索。他首先注意到 \(n!\) 等于无穷乘积 \(\prod_{k=1}^{\infty}\left(\frac{k+1}{k}\right)^{n}\frac{k}{k+n}\)，这个乘积对 n 的每个值（无论是否整数）都有定义，而且特别地，当 \(n = \frac{1}{2}\) 时由沃利斯公式它取值 \(\frac{1}{2}\sqrt{\pi}\)。这一结果与沃利斯诸结果的类比接着引导他重新拾起积分 \(\int_{0}^{1} x^{e}(1-x)^{n} dx\)（n 为整数，e 任意），它已被沃利斯考虑过。欧拉借助二项展开得到它的值 \(\frac{n!}{(e+1)(e+2)\cdots(e+n)}\)；随后一个变量代换向他表明，\(n!\) 是积分 \(\int_{0}^{1}\left(\frac{1-x^{s}}{z}\right)^{n} dx\) 当 z 趋于 0 时的极限，由此得到“第二类欧拉积分”\(n! = \int_{0}^{1}\left(\log\frac{1}{x}\right)^{n} dx\)；用同样的方法并利用沃利斯公式，他得到公式 \(\int_{0}^{1}\sqrt{\log\frac{1}{x}} dx = \frac{1}{2}\sqrt{\pi}\)。在其后的工作中，欧拉经常回到这些积分上来；他就这样发现了余元关系（{[}108 a{]}，(1)，第 XV 卷，第 82 页及第 XVII 卷，第 342 页）、公式 \(B(p,q) = \Gamma(p)\Gamma(q)/\Gamma(p+q)\)（{[}108 a{]}，(1)，第 XVII 卷，第 355 页）以及勒让德–高斯公式对应于 x = 1 的特例（{[}108 a{]}，(1)，第 XIX 卷，第 483 页）；当然，这一切都全然不关心收敛性的问题。

高斯在他关于超几何函数研究的背景下继续了对 \(\Gamma\) 函数的研究，\(\Gamma\) 函数是超几何函数的一个极限情形（{[}124 a{]}，第 III 卷，第 125–162 页）；正是在这项研究中他得到了一般的乘法公式（不久前勒让德已对 \(p = 2\) 记下了它）。其后关于 \(\Gamma\) 的工作主要涉及把这一函数推广到复域。直到最近人们才注意到，对数凸性在实域上把 \(\Gamma(x)\) 在函数方程 \(f(x + 1) = xf(x)\) 的全部解中确定到一个因子（{[}26{]}，第 149–164 页）；而阿廷已经表明{[}7 d{]}，关于 \(\Gamma(x)\) 的一切经典结果都可以简单地与这一性质联系起来。

